\section{Introduction}\label{sec:introduction}

A degeneration of the base can create slowly relaxing endomorphisms of a
Hermitian--Einstein bundle even when its connection has nonzero curvature.
The local parallel algebras, their identifications across the degenerating
regions, and the symmetries of the compact family determine which of these
endomorphisms survive globally. We study this mechanism on a compact
Calabi--Yau threefold with a nontrivial rank-four bundle. Its effective metric
variables form a five-dimensional constrained space, and the constraints
retain gluing information that the scalar dual graph of the central fibre
does not record.

Our threefolds are free quotients of a product of a fixed K3 surface and
an elliptic curve. Such Calabi--Yau threefolds are called of \emph{type K}.
Here a Calabi--Yau threefold means a smooth projective threefold with
trivial canonical bundle and $h^{1,0}=h^{2,0}=0$. The examples have infinite
fundamental group and an index-two product cover. The K\"ahler metrics used
in the analysis descend from a fixed Ricci-flat K3 metric and a smooth
ambient metric on the degenerating elliptic family. They need not be
Ricci-flat. The geometric construction, its actual bundle connection,
and these metric choices are specified in Section~\ref{sec:geometry}.

\subsection{The family and the main result}
Let $S$ be a projective K3 surface with a free Enriques involution
$\iota$. Resolve the two successive singular points of
\[
 y^2z=x(x-z)(x-t^2z)
\]
over $t=0$. This gives a projective elliptic family $\cE\to\Delta$
with central fibre a cycle of four rational curves $D_j$, indexed
modulo four. Elliptic inversion $\sigma$ acts by $j\mapsto-j$.
Set $X_t=(S\times E_t)/\langle(\iota,\sigma)\rangle$.
The central threefold has three components, with interval dual graph.
Choose the K3 metric to have volume one. The actual elliptic
component areas are $A_0,A_1,A_2,A_3=A_1>0$.

For $\theta\in\RR$ let
\[
 R=e^{i\theta\sigma_z/4},\qquad J=i\sigma_x,
 \qquad U_\theta=R^4,
 \qquad
 \sigma_x=\begin{pmatrix}0&1\\1&0\end{pmatrix},\quad
 \sigma_z=\begin{pmatrix}1&0\\0&-1\end{pmatrix}.
\]
The rank-two flat bundle $V_\theta$ on the elliptic cover has transport
$R$ at each oriented node incidence, cycle holonomy $U_\theta$,
and trivial vanishing-cycle holonomy. Its lift of $\sigma$ has square
$-I$. Tensoring this lift with $i\,d\iota$ gives a genuine
order-two action on $T^{1,0}S\otimes V_\theta$, which descends to
a rank-four bundle $\cB_{t,\theta}$. Section~\ref{sec:geometry}
constructs these transitions on an actual cover of the compact
total space and verifies the descent. In particular, the family is
not specified merely by a proposed local gluing diagram.

Let $L=\log(1/|t|)$. On
$\RR\oplus\RR^3\oplus\RR$ define the energy and mass
\begin{equation}\label{eq:intro-forms}
 \begin{split}
 q_\theta(a,X,b)&=2\bigl(\|au-X\|^2+\|X-bv_\theta\|^2\bigr),\\
 m_A(a,X,b)&=A_0a^2+2A_1\|X\|^2+A_2b^2,\\
 u&=(1,0,0),\qquad v_\theta=(\cos\theta,\sin\theta,0).
 \end{split}
\end{equation}
Write $\mu_0\le\cdots\le\mu_4$ for their generalized eigenvalues,
including zero modes. Let $\cP_2$ be the space of positive Hermitian
two-by-two matrices of determinant one, with affine invariant metric
\eqref{eq:matrix-metric}, and let $d_2$ be its geodesic distance.
The finite-dimensional configuration space is
\[
 \cF=\{(H_j)\in\cP_2^4:H_{-j}=JH_jJ^*\},
 \qquad g_M=\sum_j A_jg_{H_j}.
\]
It is a five-dimensional geodesically convex submanifold. The two
endpoint factors are one-dimensional, whereas the middle factor
has dimension three and weight $2A_1$.

The initial tuple for the effective evolution is selected from the
spatial initial metric. On each fixed normalization component, run
the HYM metric flow and take its parallel limiting metric. We call
this limit the \emph{component phase}. It is a positive metric in
the component's parallel endomorphism algebra, rather than an angular
parameter. Theorem~\ref{thm:parallel-phase} proves its existence and
quantitative dependence on the initial metric. In the degeneration,
the component restrictions admit smooth extensions across shrinking
node discs; the allowed extensions and their asymptotically vanishing
phase ambiguity are specified in Section~\ref{sec:spatial-entry}.
The norm used there controls the logarithm of the relative metric
globally in $C^0$ and controls extensions of its component-exterior
restrictions to the fixed spaces $S\times D_j$ in $C^{6,\beta}$.
It permits spatial variation in both factors and arbitrary smooth
changes of bounded amplitude inside the middle neck regions.

\begin{theorem}\label{thm:main}
For the compact family, bundle and prescribed K\"ahler metrics just
described, the following statements hold on the same smooth fibres.
\begin{enumerate}
\item The bundle has a holomorphically trivial determinant and an
explicit determinant-normalized Hermitian--Einstein metric $h_0$.
On an embedded K3 slice its second Chern number is $48$. It is
stable for $\theta\notin\pi\ZZ$ and strictly polystable for
$\theta\in\pi\ZZ$.
\item Let $\lambda_k(t,\theta)$ be the eigenvalues, indexed from
zero, of the actual endomorphism connection Laplacian
$D_\theta^*D_\theta$ on trace-free $h_0$-Hermitian endomorphisms.
There are $\varepsilon(t)\to0$ and $\gamma>0$, independent of
$\theta$, such that
\begin{equation}\label{eq:main-spectrum}
 (1-\varepsilon(t))\frac{2\pi}{L}\mu_k(\theta)
 \le\lambda_k(t,\theta)\le
 (1+\varepsilon(t))\frac{2\pi}{L}\mu_k(\theta)
 \quad(0\le k\le4),\qquad \lambda_5\ge\gamma.
\end{equation}
The kernel is one-dimensional precisely at $\theta\in\pi\ZZ$.
The comparison is relative for every positive weak mode, with no
restriction on how quickly $\theta$ approaches that locus.
\item The effective evolution is the negative gradient on $(\cF,g_M)$
of
\begin{equation}\label{eq:main-potential}
 \frac\pi L\sum_{j\bmod4}d_2(H_{j+1},RH_jR^*)^2.
\end{equation}
For every fixed $R_0>0$ and $0<\beta<1$, the nonempty open class
$\mathcal B_{R_0}(t,\theta)$ defined by the component-extension
norm \eqref{eq:extension-norm} has a selected tuple of component
phases in $\cF$. The class includes non-tensor rank-four metrics.
There is a common time $s_*(t)\to\infty$, with $s_*/L\to0$,
after which the actual equation
\[
 h^{-1}\partial_sh=-2i\Lambda_{\omega_t}F_h
\]
is approximated by the reconstructed effective evolution in global
squared $L^2$ affine distance. The error tends to zero uniformly
over this class and all $\theta$ on each interval
$[s_*,s_*+TL]$ with fixed $T$. If $\mu_0(\theta)\ge\kappa>0$,
the error tends to zero uniformly for every $s\ge s_*$.
The reconstruction and error bounds are given in
Theorems~\ref{thm:actual-tracking}, \ref{thm:all-time}
and \ref{thm:spatial-entry}.
\item Suppose the family, $L$ and the three areas are known, but
$\theta$ and the transition matrices are unknown. From the five
actual eigenvalues, the quantity
\begin{equation}\label{eq:main-inverse}
 \frac{A_0A_1^3A_2}{8}\left(\frac{L}{2\pi}\right)^5
                    \prod_{k=0}^4\lambda_k(t,\theta)
\end{equation}
recovers $\sin^2\theta$ with relative error tending uniformly to
zero when it is positive, and equals zero at reducibility. It
therefore recovers the conjugacy class of $\Ad(U_\theta)$ in
$\operatorname{SO}(3)$ uniformly. The quantitative noise bounds
are those of Theorem~\ref{thm:transport-inverse}.
\end{enumerate}
Both geometric parameters are fixed along every flow trajectory.
\end{theorem}

\begin{proof}
Propositions~\ref{prop:elliptic-family}--\ref{prop:he-stability}
prove part 1 on the stated family. Theorem~\ref{thm:five-spectrum}
gives part 2 after the exact K3 product decomposition and the
equivariant capacity comparison. Proposition~\ref{prop:finite-flow}
derives the field in part 3. Theorems~\ref{thm:actual-tracking}
and \ref{thm:all-time} prove its reconstruction and comparison with
arbitrary matching rank-four initial metrics. The fixed-bundle
phase theorem and the preparation estimate of
Theorem~\ref{thm:spatial-entry} supply such matching from the
declared open class. Theorem~\ref{thm:transport-inverse} proves
part 4. All steps use the same component areas, incidence transport
and determinant normalization. The cancellation of the rank-two
tangent factor with the degree-two quotient in
\eqref{eq:norm-cancellation} and \eqref{eq:tensor-distance}
ensures that the spectral and nonlinear coefficients agree.
\end{proof}

\subsection{Proof mechanisms and scope}
The exact product connection separates endomorphisms parallel in the
K3 direction from a complement with a fixed positive gap. The
remaining curve modes are subject to the involution. In the central
quotient, this makes the two endpoint parallel spaces lines in a
three-dimensional matrix space. Their relative angle records the
transport around the cycle upstairs, even though the downstairs
dual graph is an interval. Component averaging and neck interpolation
respect the involution. The lower spectral bound is controlled by
the section's own energy, which preserves relative accuracy when
the smallest effective eigenvalue tends arbitrarily rapidly to zero.

For the nonlinear statement, a component Green correction cancels
the leading effective velocity. Corrected endpoint geodesics and
relative-logarithm cutoffs then give a global Chern residual of
order $(1+\log L)^2/L^2$ in integral, uniformly for all physical
times. Tensor curvature and descent produce a comparison metric on
the actual rank-four bundle. An intrinsic trace inequality compares
it to exact solutions with no tensor restriction. To obtain the
open class, a separate preparation metric uses spatial component
flows. Its node profile has nonzero transverse mean curvature;
the proof controls that term using the small node volume. This
distinguishes the preparation estimate from a formal extension of
the one-dimensional equation.

The all-time assertion requires the positive-rate class in the
theorem. The weakest linear time scale near reducibility is instead
$L/\vartheta^2$, where $\vartheta$ is the distance to $\pi\ZZ$.
Section~\ref{sec:inverse-limits} proves that physical time and
degeneration generally do not commute, that rescaled time restores
commutation on a positive-rate class, and that approach to reducibility
can destroy that commutation again. For spatial initial data, a
limit at fixed positive rescaled time also requires convergence of
the selected phases. The inverse in the theorem concerns projective
bundle transport; it does not recover an integral cohomology ring,
a central lift, or a Ricci-flat metric on the threefold.

\subsection{Relation to previous work}
The smooth Enriques Calabi--Yau quotient is classical; see
Hashimoto--Kanazawa \cite{hashimoto-kanazawa} for its place among
threefolds of type K. Rank-four bundles on this very quotient
also predate the present construction. Thomas \cite[Section 2,
Theorem 2.4]{thomas} constructs polystable rank-four bundles with
trivial determinant; we use the corrected preprint for the associated
statements. Andreas--Curio \cite[Section 4, Proposition 4.2]{andreas-curio}
give a stability criterion for already non-split extensions on the
Enriques elliptic threefold, under their numerical hypotheses.
Neither the quotient geometry nor the existence of rank-four bundles
on it is the contribution claimed here. Our construction provides
a specified nonflat tangent-tensor bundle family for which the
degeneration, complete weak spectrum, actual nonlinear evolution
and projective inverse can be proved together.

Graph descriptions of small eigenvalues have a substantial history,
including Burger's relative spectral comparison for degenerating
Riemann surfaces \cite{burger}. Yang--Mills flow under degeneration
of the surface was studied by Daskalopoulos--Wentworth
\cite{daskalopoulos-wentworth}. The ambient-induced metric used
here has collapsing neck area; its zero-form mass must be computed
for that metric even though curve Dirichlet energy is conformally
invariant. Effective HYM dynamics and bundle gluing for this type
of curve degeneration are developed in
\cite{wuebben-effective,wuebben-gluing}. We give the required
equivariant estimates here and prove their relation to the actual
nonflat threefold operator and flow.

The fixed-base existence and order comparison are the classical
compact K\"ahler inputs recalled in \cite[Section 4.1]{hkkp};
the trace comparison is used in the form of
\cite[Proposition 2.4]{wang}. These inputs are valid in arbitrary
complex dimension. They do not supply uniform derivative or heat
kernel estimates for our degenerating family. The relative weak
spectral comparison, the global residual, the spatial preparation
and their parameter/time domains are proved below.
